% Mathematics agent LaTeX fragment for BioSci manuscript
% Drop into biological_science/workspace/manuscript/sections/02_results_framework.tex and 07_methods.tex
% Notation follows DS's 00_design_lrs_framework.md v0.1
% Author: Mathematics Agent (agent-mozusggu), 2026-05-10

% ============================================================
% For sections/02_results_framework.tex
% ============================================================

\paragraph{Theoretical guarantees for LRS.}
We establish three theorems that jointly anchor the LRS framework (see Supplementary Section~S1 for full proofs).

\textbf{Theorem 1 (minimax floor of label-free detection).}
Let the observed cell population follow a mixture $p = (1-\pi)\, p_{\text{abund}} + \pi\, p_{\text{rare}}$ with Wasserstein separation $W_2(p_{\text{rare}}, p_{\text{abund}}) \geq \Delta$, and let $\psi$ be any label-free test with Type-I error~$\alpha$. There exists an absolute constant $c_1 > 0$ such that if
\begin{equation}
  n\,\pi^{2}\,\Delta^{2}/(\sigma^{2} d)
  \;\leq\; c_{1}\,\log \bigl(1/(\alpha\beta)\bigr),
  \label{eq:minimax-floor}
\end{equation}
no label-free test achieves power at least $1-\beta$. For the operating regime $\pi\in[10^{-4},10^{-2}]$, $\Delta\in[0.3,1.5]\,\sigma$, $d=32$, $\alpha=\beta=0.05$, inequality~\eqref{eq:minimax-floor} gives $n \gtrsim 767/(\pi^{2}\Delta^{2})$, matching the scalability claim at the Kang~2024 atlas scale ($n\approx 4.9\times 10^{6}$).

\textbf{Theorem 2 (topological stability of $P_t$).}
If the integrator induces a per-cell displacement $\|\phi(x)-x\|_{\infty}\leq \varepsilon$, then the topological-persistence score satisfies
\begin{equation}
  \bigl|\,P_t(w)-1\,\bigr| \;\leq\; C(\rho_{\max},\operatorname{diam}(\Omega),d)\cdot \varepsilon / L_{\text{pre}}(w),
\end{equation}
where $L_{\text{pre}}(w)$ is the longest $H_0$ bar of the pre-integration witness-complex filtration on the seed. Consequently $P_t(w)<0.4$ is not explainable by mild integration smoothing and constitutes strong evidence of topological collapse. The proof adapts Cohen-Steiner--Edelsbrunner--Harer (2007) to the DTM-filtered witness complex (Chazal--Cohen-Steiner--M\'erigot 2011).

\textbf{Theorem 3 (cross-batch witness identifiability).}
Under the transcriptomic-separation condition (A1), the HVG-space common-factor condition (A2), the kNN-density consistency condition (A3), and the common-factor batch-independence condition (A4) (Methods~\S2.3), for every pair of batches there exists a constant $c_3(d, \sigma, \Delta, \rho_{\max})$ such that the true rare population is recovered as a witnessed seed $w\in W$ with probability at least $1-\delta$ whenever
\begin{equation}
  n_b\,\pi_b^{2} \;\geq\; c_{3}\,\log(1/\delta)/\gamma^{2}
  \quad\text{for all } b \in \mathcal{B}.
  \label{eq:identifiability}
\end{equation}
With $B\geq 10$ batches and pooling, the effective sample-size condition becomes $n\pi^{2}\geq B\,c_3\log(1/\delta)/\gamma^{2}$, which at Kang-atlas scale permits $\pi\geq 6\times 10^{-4}$. This defines the LRS detectable-regime envelope (Fig.~1d).

Theorems~1 and~3 together establish a matched upper and lower bound on label-free detectability: LRS saturates the information-theoretic floor, up to constants.

% ============================================================
% For sections/07_methods.tex, subsection "Mathematical formalization"
% ============================================================

\subsection*{Mathematical formalization of LRS}

\paragraph{Setting.}
Observations are modelled as $n$ independent draws from a mixture $p = (1-\pi)p_{\text{abund}} + \pi p_{\text{rare}}$ on $\mathbb{R}^{d}$ ($d=32$ PCA-reduced HVG space by default). An integration is any measurable map $T:\bigsqcup_{b}\mathbb{R}^{n_{b}\times d}\to \mathbb{R}^{n\times d}$ acting cell-wise. Preservation of a biological component $C\subseteq \operatorname{supp}(\mu_{\text{bio}})$ requires both mass matching $|\mu_{\text{post}}(\widehat{C}) - \pi(C)|\leq \eta\cdot \pi(C)$ and Hausdorff locality $d_{H}(\widehat{C}, T_{\#}\mu_{\text{bio}}|_{C})\leq \eta\cdot\operatorname{diam}(T_{\#}\mu_{\text{bio}}|_{C})$, for some tolerance $\eta$.

\paragraph{Blind spot of label-based evaluation (Theorem 0).}
Any evaluation functional $E$ that depends only on $(T,y)$ where $y$ denotes annotator labels is invariant under local integration perturbations supported on an unannotated component $C$. Concretely, for any such $E$ and any integration $T_{0}$, there exists $T_{1}$ with $E(T_{1},y)=E(T_{0},y)$ for every $y$ yet with $T_{1}$ crushing $C$ into an adjacent abundant component. Hence label-based metrics (ARI, NMI, ASW, graph-connectivity, iLISI, cLISI, kBET, scIB-E extensions) are blind by construction to destruction of unannotated rare populations (Supplementary Section~S1, Theorem~S1).

\paragraph{Assumptions of LRS identifiability (Theorem 3).}
(A1) Transcriptomic separation: $W_{2}(p_{\text{rare}}, p_{\text{abund}})\geq \Delta$.
(A2) HVG-space common factor: the HVG projection of $p_{\text{rare}}$ has a batch-invariant mean satisfying $\cos(\mu_{\text{rare}}^{(b)}, \mu_{\text{rare}}^{(b')})\geq \sigma_{\text{thresh}} + \gamma$ for all pairs $b,b'$.
(A3) kNN density consistency: the per-batch kNN density estimator $\rho_{b}$ with parameters $(k,K)$ concentrates with sub-Gaussian tails at rate $\eta = O(\sqrt{\log(n_{b})/(n_{b} k \Delta^{d})})$.
(A4) Common-factor independence: batch labels are conditionally independent of the rare/abundant label given HVG coordinates.

\paragraph{Null-distribution calibrations.}
LRS($w$) is calibrated under three label-free nulls: (N1) batch-label permutation; (N2) within-cell gene-shuffle; (N3) per-batch HVG rotation. Under each, LRS($w$) is asymptotically degenerate at~1 up to $O(1/\sqrt{n})$ fluctuations. Bootstrap 95\% CIs are obtained by resampling $w\in W$ with replacement ($R=1000$); the delta-method variance $\operatorname{Var}(\mathrm{LRS}(f))=O(1/|W|)$ gives closed-form sample-size estimates. For atlases with potential batch-hidden confounding of rare populations (e.g., Kang~2024 where patient metadata is non-random across batches), report $\max(p_{\text{N1}}, p_{\text{N3}})$ as the confounding-robust p-value.

\paragraph{REAL-weighted scVI identifiability (Theorem 5).}
The REAL penalty $R(z,W)=\sum_{w\in W}\|\operatorname{center}_{q}(z\mid C_{w}) - \operatorname{center}_{p_{\text{prior}}}(z\mid C_{w})\|^{2}$ added to the scVI ELBO with coefficient $\lambda\geq 0$ preserves scVI's identifiability up to the native batch-equivariant rotation group. The MAP point estimate on $C_{w}$ is biased toward $\operatorname{center}_{p_{\text{prior}}}(z\mid C_{w})$ by $O(\lambda \sigma_{z}^{2})$; this is an inductive bias, not a distortion of identifiability.
